\section{Sampling Distribution of the Mean and Unbiased Variance}

The sampling distribution of the mean is fundamental for statistical inference. Its detailed treatment is carried out in Section \ref{sec:3.4}; here we summarize the main results.

If $X_1, X_2, \ldots, X_n$ are independent and identically distributed (iid) random variables with mean $\mu$ and variance $\sigma^2$, and $\bar{X} = \frac{1}{n}\sum X_i$, then
\begin{align}
 \label{eq:3.2.7}
 E(\bar{X}) &= \mu, \\
 \label{eq:3.2.8}
 \text{Var}(\bar{X}) &= \frac{\sigma^2}{n}, \\
 \label{eq:3.2.9}
 \text{SD}(\bar{X}) &= \frac{\sigma}{\sqrt{n}}.
\end{align}

Furthermore, by the Central Limit Theorem, when $n$ is sufficiently large,
\begin{align}
 \label{eq:3.2.10}
 \frac{\bar{X} - \mu}{\sigma/\sqrt{n}} \xrightarrow{d} N(0,1).
\end{align}

\begin{observacion}
The exact distribution of $\bar{X}$ is \emph{normal} if the original population is normal, and approximately normal for large $n$ in other cases. This difference is crucial: in the first case we speak of an exact $Z$; in the second, of an approximate $Z$.
\end{observacion}



\subsection{Statistics and the Unbiased Sample Variance}

\begin{definicion}[Statistic]
Let $X_1, X_2, \ldots, X_n$ be a simple random sample (iid random variables) from a population. A \emph{statistic} is any function $T = g(X_1, \ldots, X_n)$ of the sample that does not depend on unknown population parameters. Every statistic is itself a random variable with its own probability distribution, called the \emph{sampling distribution} of $T$.
\end{definicion}

The sample mean $\bar{X}$ is the most commonly used statistic; the second most important is the \emph{sample variance}.

\begin{definicion}[Sample variance]
 \label{eq:5.1.1}
Given a simple random sample $X_1, \ldots, X_n$ with sample mean $\bar{X}$, the \emph{sample variance} is defined as
\begin{align}
 S^2 = \frac{1}{n-1}\sum_{i=1}^n (X_i - \bar{X})^2.
\end{align}
\end{definicion}

\begin{observacion}[Why divide by $n-1$ instead of $n$?]
The divisor $n-1$ (known as \emph{Bessel's correction}) is not arbitrary: it is precisely what makes $S^2$ an \emph{unbiased estimator} of the population variance $\sigma^2$, as shown below.
\end{observacion}

\begin{teorema}[Unbiasedness of the sample variance]
 \label{eq:5.1.2}
If $X_1, \ldots, X_n$ are iid with mean $\mu$ and variance $\sigma^2 < \infty$, then
\begin{align}
 E(S^2) = \sigma^2.
\end{align}
\end{teorema}

\begin{proof}
We write the sum of squares centered at $\bar{X}$ in terms of the population mean $\mu$, adding and subtracting $\mu$:
\begin{align}
 \sum_{i=1}^n (X_i - \bar{X})^2
 &= \sum_{i=1}^n \left[(X_i - \mu) - (\bar{X} - \mu)\right]^2 \\
 &= \sum_{i=1}^n (X_i - \mu)^2 - 2(\bar{X}-\mu)\sum_{i=1}^n(X_i-\mu) + n(\bar{X}-\mu)^2 \\
 &= \sum_{i=1}^n (X_i - \mu)^2 - n(\bar{X}-\mu)^2,
\end{align}
where we used $\sum_{i=1}^n (X_i - \mu) = n(\bar{X}-\mu)$. Taking expectations and using $E[(X_i-\mu)^2] = \sigma^2$ for each $i$, and $E[(\bar{X}-\mu)^2] = \text{Var}(\bar{X}) = \sigma^2/n$:
\begin{align}
 E\left[\sum_{i=1}^n (X_i - \bar{X})^2\right] = n\sigma^2 - n\cdot\frac{\sigma^2}{n}
 = (n-1)\sigma^2.
\end{align}
Dividing by $n-1$:
\begin{align}
 E(S^2) = E\left[\frac{1}{n-1}\sum_{i=1}^n (X_i - \bar{X})^2\right]
 = \frac{(n-1)\sigma^2}{n-1} = \sigma^2. \qedhere
\end{align}
\end{proof}

\begin{observacion}[The ``natural'' estimator is biased]
The estimator that divides by $n$ instead of $n-1$, $\hat{\sigma}^2 = \frac{1}{n}\sum(X_i-\bar{X})^2 = \frac{n-1}{n}S^2$, systematically \emph{underestimates} the population variance:
\begin{align}
 E(\hat{\sigma}^2) = \frac{n-1}{n}\sigma^2 < \sigma^2.
\end{align}
Intuitively, $\bar{X}$ is, by construction, closer to each $X_i$ than the unknown population mean $\mu$, so the sum of squares around $\bar{X}$ underestimates the true spread; dividing by $n-1$ instead of $n$ (``losing one degree of freedom'' from having estimated $\mu$ with $\bar{X}$) corrects exactly for that bias.
\end{observacion}

\begin{ejemplo}
 \label{exmp:sample-mean-and-unbiased-variance}
A simple random sample of size $n=5$ yields the observations $\{12, 15, 11, 18, 14\}$. Compute the sample mean $\bar{X}$ and the unbiased sample variance $S^2$.
\end{ejemplo}

\begin{solucion}
The sample mean is
\begin{align}
 \bar{X} = \frac{12+15+11+18+14}{5} = \frac{70}{5} = 14.
\end{align}
The squared deviations are $(12-14)^2=4$, $(15-14)^2=1$, $(11-14)^2=9$, $(18-14)^2=16$, $(14-14)^2=0$, summing to $30$. Therefore,
\begin{align}
 S^2 = \frac{1}{n-1}\sum_{i=1}^n (X_i-\bar{X})^2 = \frac{30}{4} = 7.5.
\end{align}
\end{solucion}



\subsection{Central Limit Theorem: Asymptotic Convergence}

\begin{definicion}[Convergence in distribution]
A sequence of random variables $Y_n$ \emph{converges in distribution} to $Y$ (written $Y_n \xrightarrow{d} Y$) if $F_{Y_n}(y) \to F_Y(y)$ at every continuity point of $F_Y$.
\end{definicion}

\begin{teorema}[Central Limit Theorem (Lindeberg--L\'evy)]
 \label{eq:5.2.1}
Let $X_1, X_2, \ldots$ be iid random variables with mean $\mu$ and variance $\sigma^2 < \infty$. Then
\begin{align}
 Z_n = \frac{\bar{X}_n - \mu}{\sigma/\sqrt{n}} \xrightarrow{d} N(0,1)
 \quad \text{as } n \to \infty.
\end{align}
\end{teorema}

\begin{proof}[Proof sketch via MGF]
Suppose $M_{X_1}(t)$ exists in a neighborhood of $t=0$. The MGF of $Z_n$ is
\begin{align}
 M_{Z_n}(t) = e^{-\mu t \sqrt{n}/\sigma}\left[M_{X_1}\!\left(\frac{t}{\sigma\sqrt{n}}\right)\right]^n.
\end{align}
Expanding $\log M_{X_1}(s)$ in a Taylor series around $s=0$:
\begin{align}
 \log M_{X_1}(s) = \mu s + \frac{\sigma^2 s^2}{2} + O(s^3).
\end{align}
Substituting $s = t/(\sigma\sqrt{n})$ and multiplying by $n$:
\begin{align}
 \log M_{Z_n}(t) = n\left[\frac{\mu t}{\sigma\sqrt{n}} + \frac{t^2}{2n} + O(n^{-3/2})\right]
 - \frac{\mu t\sqrt{n}}{\sigma} = \frac{t^2}{2} + O(n^{-1/2}) \;\longrightarrow\; \frac{t^2}{2}.
\end{align}
Hence $M_{Z_n}(t) \to e^{t^2/2}$, the MGF of $N(0,1)$; by L\'evy's uniqueness and continuity theorem, $Z_n \xrightarrow{d} N(0,1)$.
\end{proof}

\begin{observacion}[Independence from the original distribution]
The result does not depend on the shape of the distribution of $X_i$ (discrete, continuous, symmetric, skewed) --- it only requires finite mean and variance. This universality is why the CLT is the cornerstone of classical statistical inference.
\end{observacion}

\subsubsection{Rate of convergence: the Berry--Esseen theorem}

The CLT guarantees convergence as $n \to \infty$, but says nothing about how fast. The following result quantifies the error of the normal approximation for finite $n$.

\begin{teorema}[Berry--Esseen]
 \label{eq:5.2.2}
Let $X_1,\ldots,X_n$ be iid with $E|X_i-\mu|^3 = \rho < \infty$. There exists a universal constant $C$ (known to satisfy $C < 0.5$) such that
\begin{align}
 \sup_z \left|F_{Z_n}(z) - \Phi(z)\right| \le \frac{C\rho}{\sigma^3\sqrt{n}}.
\end{align}
\end{teorema}

\begin{observacion}[The practical rule $n \ge 30$]
The Berry--Esseen theorem confirms that the error of the normal approximation decays as $O(1/\sqrt{n})$: quadrupling $n$ halves the maximum error. The popular ``$n \ge 30$'' rule is a reasonable heuristic for moderately symmetric populations with finite moments, but heavily skewed or heavy-tailed distributions (i.e., with large $\rho$) require larger $n$ to reach the same precision.
\end{observacion}

\begin{ejemplo}
 \label{exmp:5.2.1}
Service time at a call center follows an exponential distribution with mean $\mu = 1/\lambda = 4$ minutes and standard deviation $\sigma = 4$. A sample of $n=64$ calls is taken. What is the approximate probability that the average service time exceeds $4.5$ minutes?
\end{ejemplo}

\begin{solucion}
By the CLT, $\bar{X}_{64} \approx N(4, \sigma^2/n) = N(4, 16/64) = N(4, 0.25)$, with standard deviation $\text{SD}(\bar{X}) = 0.5$.
\begin{align}
 P(\bar{X} > 4.5) \approx P\!\left(Z > \frac{4.5-4}{0.5}\right) = P(Z>1) = 1-\Phi(1)
 \approx 0.1587.
\end{align}
\end{solucion}
